% RESEARCH_DRAFT_SHA256: 61248ec91f196f7f67c46395d7a32ee6e3a23fb8842c14d6de48f504bd4e0e04
% Layout and source-reference wrappers expanded for the web renderer.
\documentclass[11pt]{article}
\usepackage[margin=25mm]{geometry}
\usepackage{fontspec,unicode-math}
\setmainfont{Latin Modern Roman}
\setmathfont{Latin Modern Math}
\usepackage{amsmath,amssymb,mathtools,xurl,hyperref}
\hypersetup{colorlinks=true,urlcolor=blue,linkcolor=blue}
\setcounter{secnumdepth}{0}
\setlength{\parindent}{0pt}
\setlength{\parskip}{5pt}
\setlength{\emergencystretch}{3em}
\newcommand{\sref}[2]{\hyperlink{source:#1:#2}{[S#2]}}
\newcommand{\cataloguescope}[1]{\paragraph{Source record.}#1}
\newcommand{\floor}[1]{\left\lfloor#1\right\rfloor}
\begin{document}
\section*{Research-attempt provenance}
This standalone research attempt retains the complete original definition body below,
including its historical source and preparation notes. Those notes describe the
earlier catalogue preparation; the new attempt and its verification are identified
separately after that body. No original ID, catalogue status, or source statement is changed.

\section{An ancestral-observation certificate for likelihood unimodality}\label{statement}

\subsection{Definitions and mathematical statement}

\paragraph{Shared notation.}Standard number systems are $\mathbb N,\mathbb Z,\mathbb Q,\mathbb R,\mathbb C$. Specialized notation and index ranges follow the cited source and the definitions below.


In a graphical model, observed variables enter the data while hidden variables are summed or integrated out. The likelihood is the probability or density of the observed sample as a function of model parameters. Unimodality concerns a single mode under specified parameterization and boundary conventions; hidden variables can create nonidentifiability and multiple likelihood maxima.
\paragraph{Record scope and qualifications.}
The catalogue title is retained as metadata. Named certificates, contracts, bounds or frameworks in that title are not asserted as results of the original source extract.
The graph alone does not specify distributions, cardinalities, parameterization, or quantification over samples. Its source question is retained without claiming a graph-only criterion for unimodality. The supplied catalogue marks this record partially solved; its present research status was not independently checked. Its supplied reconstruction is labelled unverified.
The source text below is rendered literally where necessary; damaged formulas and composite extracts are not silently repaired into a new mathematical claim.
The supplied catalogue marks this formulation as an unverified reconstruction; missing assumptions are not certified by the source link.
\paragraph{Statement recorded in the supplied source.}
Question 40 (Wynn). Which nodes being hidden imply the multimodality/ unimodality of the likelihood? Given a particular graph, what nodes must be observed to ensure unimodal-ity? \href{https://aimath.org/WWN/compalgstat/compalgstat.pdf}{[S2]}
\paragraph{Source status.}
The supplied catalogue labels this source record partially solved; that metadata was checked against the supplied snapshot on 2026-10-11. The label is preserved for provenance and is not an independent verification of present open status or mathematical correctness.
\subsection{Short English statement}
The source asks how choices of hidden and observed nodes affect likelihood modality.
\subsection{Sources}
\begin{itemize}
\item[S1] Ulam AI and the UnsolvedMath Contributors, UnsolvedMath: A Curated Collection of Open Mathematics Problems, catalogue entry. Catalogue curation and metadata: CC BY 4.0. The source extract is retained; definitions, English summaries and qualifications are editorial additions.
\url{https://huggingface.co/datasets/ulamai/UnsolvedMath}.
\item[S2] AIM Workshop Problem Lists, source linked by the catalogue.
\url{https://aimath.org/WWN/compalgstat/compalgstat.pdf}.
\end{itemize}
\paragraph{Preparation and review.}
The original source and UnsolvedMath contributors retain attribution for the source extract. Codex (GPT-6; exact serving revision not exposed) assisted with definitions, English summaries, scope notes and formatting. No new human mathematical authorship or independently verified solution is asserted. Catalogue identity checked against the complete supplied registry and current repository. Distinct upstream records, including repeated questions and broad or damaged extracts, are retained for complete coverage of the source CS category. A repeated record is not a new independent conjecture. Complete mathematical-equivalence and current-literature review has not been performed. Preparation and final review follow contribution prompts C and D, with source gaps disclosed. Structural and rendering checks do not establish a new theorem or open status.
\label{end:statement}

\section{Literature and context check}
The original AIM Computational Algebraic Statistics
 workshop PDF was consulted. The preserved extract does not specify every model,
 parameterization or performance criterion needed by its broad question. The
 following declared model is an actual test of one route, not a replacement
 target or a claim of novelty for elementary statistical or algebraic facts.
\section{Research attempt: Even fully observed likelihood modality depends on parameterization}
For a fully observed finite discrete DAG with unrestricted conditional
probability rows, the likelihood factors over rows as
$\prod_{i,a,k}\theta_{i,k\mid a}^{n_{i,a,k}}$.
Its log is a sum of nonnegative counts times logarithms, so it is concave
on the product of row simplices: $\log$ has negative second derivative, and
nonnegative sums preserve concavity. A positive-total row is maximized by
its empirical frequencies. To prove this including zero counts, write the
positive empirical probabilities as $r_k$. The inequality
$\log t\le t-1$ gives
$\sum r_k\log(\theta_k/r_k)\le\sum_{r_k>0}\theta_k-1\le0$.
Equality requires the stated row and zero mass on zero-count cells.
A row with no observations is entirely flat. Thus all local maxima in this
convex parameter space are global, while uniqueness can fail through flat rows.

This does not yield a graph-only criterion. A graph with a single observed
Bernoulli node and no hidden nodes has a unique fitted success probability
$p=n_1/N$ for mixed data. Reparameterize it by $p=\theta^2$,
$\theta\in[-1,1]$. The same model and data have two distinct parameter maxima
$\pm\sqrt{n_1/N}$; the graph and observed-node set have not changed. If one
counts modes of observable distributions, both maxima represent one mode;
if one counts parameter modes they are distinct. The question must fix which
notion it intends.

Even natural unrestricted rows do not ensure a unique maximum for every
sample: if a parent configuration is absent, its whole conditional row is
unconstrained by the likelihood. Concavity does not remove this nonidentifiability.
Therefore a proposed answer based solely on whether all nodes are observed
overstates what complete-data factorization proves. Hidden-node placement
requires still more information about state sizes, restrictions and sample
quantifiers before a precise classification can be attempted.
\section{Verification and adversarial audit}
The concavity and row-maximizer argument handles zero counts and simplex
boundaries. The single-node example uses mixed data to avoid boundary ambiguity.
Parameter modality, quotient modality and flat nonidentifiability are kept
separate.
\paragraph{Audit corrections.} Concavity was not equated with unique attainment, and hidden-node placement alone was not asserted to determine modality.
\paragraph{External computation.} None; symbolic arguments and boundary checks only.
\section{Scope and remaining gaps}
The source does not choose a parameterization, distribution family or
meaning of mode. The first gap is a precise target invariant under the intended
equivalence, followed by a graph-and-model theorem beyond the elementary
complete-data convex case.
\paragraph{Integration review.} Codex (GPT-6; exact serving revision not exposed) reviewed the written arguments and expanded layout and source-reference wrappers. Repository metadata and the production web renderer were checked using Python 3.13 and JavaScript with MathJax 3.2.2. This model-assisted review is not independent human proof verification; no PDF compilation or formal proof checking is asserted.

\section{Completion Estimate}
\noindent\textbf{COMPLETION: 0\%}
The calculations below the preserved source statement investigate a concrete route,
but establish no substantiated new progress toward the entire original target.
Known special cases and failed sufficient conditions are not credited as a solution.
\section{References}
The original source references and catalogue attribution are preserved in the
Sources section above. The following URLs were actually consulted for this attempt;
retention of an additional original reference does not claim fresh consultation of it.
\begin{itemize}
\item \url{https://aimath.org/WWN/compalgstat/compalgstat.pdf} (consulted 2026-10-11). 
\end{itemize}
\section{Model and contributor disclosure}
OpenAI gpt-6-astra, requested reasoning configuration ultra, generated this research
attempt on 2026-10-11. The exact serving revision is not exposed. Codex GPT-6
performed orchestration, repository-format review, and a model-assisted mathematical
audit of the written arguments; its exact serving revision
is also not exposed. The mathematical route, boundary checks and final adversarial
self-audit followed contribution prompts A then B in
\texttt{docs/contribution\_prompts.md}, after reading \texttt{CONTRIBUTING.md}.
Python 3.14.2 assembled the standalone documents; any mathematical computation is
specified in Verification. No theorem prover was used. No human mathematical
authorship, substantial human proof correction, or independent human proof checking
is asserted. Model-assisted review, source retrieval, and structural checks do not
establish correctness or novelty. Full private reasoning transcripts are unavailable
as reproducibility artifacts; the complete arguments, inputs and outcomes actually
used are stated here or in the companion batch audit.
\section{Final status}
Unresolved as a new research contribution. This is a concrete route investigation,
not a claim of new complete resolution or a current-status certification. A known
answer reproduced for context does not become a new contribution. The catalogue's
source-status label is preserved; this attempt status does not assert that a
source-labelled solved question is open.
\end{document}
